% Options for packages loaded elsewhere
\PassOptionsToPackage{unicode}{hyperref}
\PassOptionsToPackage{hyphens}{url}
\documentclass[
]{article}
\usepackage{xcolor}
\usepackage{amsmath,amssymb}
\setcounter{secnumdepth}{-\maxdimen} % remove section numbering
\usepackage{iftex}
\ifPDFTeX
  \usepackage[T1]{fontenc}
  \usepackage[utf8]{inputenc}
  \usepackage{textcomp} % provide euro and other symbols
\else % if luatex or xetex
  \usepackage{unicode-math} % this also loads fontspec
  \defaultfontfeatures{Scale=MatchLowercase}
  \defaultfontfeatures[\rmfamily]{Ligatures=TeX,Scale=1}
\fi
\usepackage{lmodern}
\ifPDFTeX\else
  % xetex/luatex font selection
\fi
% Use upquote if available, for straight quotes in verbatim environments
\IfFileExists{upquote.sty}{\usepackage{upquote}}{}
\IfFileExists{microtype.sty}{% use microtype if available
  \usepackage[]{microtype}
  \UseMicrotypeSet[protrusion]{basicmath} % disable protrusion for tt fonts
}{}
\makeatletter
\@ifundefined{KOMAClassName}{% if non-KOMA class
  \IfFileExists{parskip.sty}{%
    \usepackage{parskip}
  }{% else
    \setlength{\parindent}{0pt}
    \setlength{\parskip}{6pt plus 2pt minus 1pt}}
}{% if KOMA class
  \KOMAoptions{parskip=half}}
\makeatother
\usepackage{longtable,booktabs,array}
\usepackage{caption}
\captionsetup[table]{skip=6pt}
\newcounter{none} % for unnumbered tables
\usepackage{calc} % for calculating minipage widths
% Correct order of tables after \paragraph or \subparagraph
\usepackage{etoolbox}
\makeatletter
\patchcmd\longtable{\par}{\if@noskipsec\mbox{}\fi\par}{}{}
\makeatother
% Allow footnotes in longtable head/foot
\IfFileExists{footnotehyper.sty}{\usepackage{footnotehyper}}{\usepackage{footnote}}
\makesavenoteenv{longtable}
\setlength{\emergencystretch}{3em} % prevent overfull lines
\providecommand{\tightlist}{%
  \setlength{\itemsep}{0pt}\setlength{\parskip}{0pt}}
\usepackage{bookmark}
\IfFileExists{xurl.sty}{\usepackage{xurl}}{} % add URL line breaks if available
\urlstyle{same}
% fallback for those not using the hyperref driver hyperxmp:
\makeatletter
\@ifundefined{xmpquote}{\newcommand{\xmpquote}[1]{#1}}{}
\makeatother
\hypersetup{
  pdftitle={Iterative Symbolic Refit Matches Bayesian Experimental Design},
  hidelinks,
  pdfcreator={LaTeX via pandoc}}

\title{Iterative Symbolic Refit Matches Bayesian Experimental Design}
\author{}
\date{}

\begin{document}
\maketitle

\section{Iterative Symbolic Refit Matches Bayesian Experimental Design
for Physical Law
Discovery}\label{iterative-symbolic-refit-matches-bayesian-experimental-design-for-physical-law-discovery}

\textbf{11-qlab}

\emph{Preprint, November 2026}

\begin{center}\rule{0.5\linewidth}{0.5pt}\end{center}

\subsection{Abstract}\label{abstract}

We benchmark three discovery strategies on 23 physical laws spanning
classical mechanics, thermodynamics, electromagnetism,
quantum-gravity-adjacent hypothetical extensions, and interferometric
phase forms at 2\% multiplicative noise. Initial convergence checking
based on symbolic equality was found to reject noise-fitted true forms;
we replace it with a numeric-ratio check sampled at data-driven
parameter corners, which raises measured convergence rates from 48\% to
70\% for single-shot discovery and from 74\% to 91\% for iterative
refit. Under the corrected check, \textbf{iterative refit converges on
21/23 laws using 5-point initial samples at 8.8 seconds median},
dominating single-shot discovery (16/23, 15 points, 13.5s) on
convergence, data efficiency, and wall time simultaneously. Bayesian
optimal experimental design (BOED) converges on all 8 single-variable
laws it can attempt, but its 1-D proposal framework provides
structurally zero coverage of the 15 multi-variable laws in the
benchmark. Two residual failures --- \texttt{cow\_phase} (3-variable
product with a squared term) and \texttt{atom\_grav\_gradient}
(3-variable two-term offset) --- identify concrete directions for future
work. We report these failures as diagnostic measurements rather than
suppressing them, and we include the initial incorrect measurement
alongside the corrected one as part of the reproducibility record.

\textbf{Code and benchmark:} https://github.com/11-qlab/pd

\begin{center}\rule{0.5\linewidth}{0.5pt}\end{center}

\subsection{1. Introduction}\label{introduction}

Symbolic regression --- the problem of recovering a closed-form
expression from numerical data --- has advanced rapidly in the last
decade {[}Schmidt \& Lipson 2009; Cranmer 2023; Udrescu \& Tegmark
2020{]}. Systems such as PySR and AI Feynman can recover canonical
physics formulas from noisy data in seconds, and the recent surge of
LLM-driven scientific agents has extended the scope from formula
recovery to closed-loop experimental design {[}Boiko et al.~2023;
Rappaport et al.~2024; Lu et al.~2024{]}.

A natural question follows: \textbf{does adaptive experimental design
improve symbolic law discovery?} Concretely: given a set of candidate
laws and a budget of measurements, does it help to iterate --- proposing
each new measurement based on current candidates --- or is a single
well-designed batch sufficient?

This paper answers that question empirically. We compare three discovery
strategies on 23 physical laws:

\begin{enumerate}
\def\labelenumi{\arabic{enumi}.}
\tightlist
\item
  \textbf{Single-shot} --- one batch of measurements, one discovery
  pass.
\item
  \textbf{Random growth} --- start with a small batch, add one random
  measurement per iteration, re-discover.
\item
  \textbf{Closed-loop BOED} --- start with a small batch, propose the
  next measurement via Bayesian optimal experimental design,
  re-discover.
\end{enumerate}

Our main findings, each backed by measurements:

\textbf{Finding 1.} Iterative random refit (\texttt{random\_growth})
strictly dominates single-shot discovery: 21/23 vs 16/23 convergence at
1/3 the data and 1.5× the speed.

\textbf{Finding 2.} BOED converges on all 8 single-variable laws it can
attempt, but the 1-D proposal framework is structurally inapplicable to
15 of 23 laws. On the laws it \emph{can} attempt, BOED matches random
refit at 5 points --- providing no measurable benefit for the tested
class.

\textbf{Finding 3.} Multi-variable forms are solvable by iterative
refit. Our earlier claim {[}internal report{]} that multi-variable
search was fundamentally harder than 1-variable search was an artifact
of test-range design, not a fundamental search limit.

We also report a methodological correction: our initial convergence
check, based on symbolic equality, rejected noise-fitted true forms.
Replacing it with a numeric-ratio check raises measured convergence
across both primary methods without changing the underlying search. We
report both measurements in the interest of reproducibility.

\begin{center}\rule{0.5\linewidth}{0.5pt}\end{center}

\subsection{2. Related Work}\label{related-work}

\textbf{Symbolic regression engines.} PySR {[}Cranmer 2023{]} uses
genetic programming with distributed search. AI Feynman {[}Udrescu \&
Tegmark 2020{]} applies physics-specific priors (dimensional analysis,
symmetry) to decompose expressions. Both focus on the search problem;
neither closes the loop with experiment design.

\textbf{Automated scientific discovery platforms.} Quine {[}Rappaport et
al.~2024{]} closes the loop between hypothesis generation, experiment
design, and bio-lab execution. Coscientist {[}Boiko et al.~2023{]} uses
LLMs to drive chemistry experiments. The Sakana AI Scientist {[}Lu et
al.~2024{]} generates ML research papers end-to-end. None of these
benchmarks symbolic law discovery across physical domains.

\textbf{Bayesian optimal experimental design.} BOED has a long history
in the experimental design literature {[}Chaloner \& Verdinelli 1995;
Rainforth et al.~2024{]}. Its application to symbolic regression has
been limited. Our benchmark is, to our knowledge, the first systematic
comparison of BOED-style adaptive sampling against non-adaptive
baselines for physical law recovery.

\textbf{Physics-specific regression.} AI Feynman applies dimensional
analysis and symmetry constraints during search. Our benchmark laws span
fundamental and derived forms; we do not use physical priors during
search beyond basic search-space constraints.

\begin{center}\rule{0.5\linewidth}{0.5pt}\end{center}

\subsection{3. Methods}\label{methods}

\subsubsection{3.1 Search pipeline}\label{search-pipeline}

Our pipeline (\texttt{purediscover.discover\_and\_classify}) applies
four search strategies in sequence, stopping when the first returns a
candidate whose fit error on the training data is below tolerance:

\begin{enumerate}
\def\labelenumi{\arabic{enumi}.}
\tightlist
\item
  \textbf{\texttt{brute\_search}} --- beam search over a small grammar
  (up to 10 nodes).
\item
  \textbf{\texttt{search\_exponential}} --- log-linear regression over a
  physics-motivated feature basis (features: \texttt{x}, \texttt{1/x},
  \texttt{x²}, \texttt{log(x)}, \texttt{sqrt(x)}; max 2 features in the
  exponent).
\item
  \textbf{\texttt{search\_offset}} --- 2-parameter Laurent fit
  \texttt{c0\ +\ c1/x} for offset forms.
\item
  \textbf{\texttt{search\_composition}} --- small-candidate enumeration
  (B candidates), data peel, independent A search.
\end{enumerate}

All strategies use the same underlying tolerance and convergence check.

\subsubsection{3.2 Benchmark methods}\label{benchmark-methods}

\textbf{\texttt{single\_shot}} --- generate 15 measurements once, run
discovery once, check convergence.

\textbf{\texttt{random\_growth}} --- generate 5 measurements, run
discovery, check convergence. If not converged, add one random
measurement, re-discover, repeat up to 10 iterations.

\textbf{\texttt{closed\_loop\_boed}} --- same as random\_growth but the
additional measurement is proposed by Bayesian optimal experimental
design. Our BOED implementation is single-variable only: it selects the
x value maximizing between-candidate variance weighted by
unexplored-region coverage.

\subsubsection{3.3 Convergence check}\label{convergence-check}

A fitted expression \texttt{found} is converged to \texttt{truth} iff
\texttt{found} is a nonzero scalar multiple of \texttt{truth}. Two
versions were tested:

\textbf{Symbolic (initial).}
\texttt{sp.simplify(found\ -\ truth)\ ==\ 0}, or
\texttt{sp.simplify(found\ /\ truth)} has no free symbols.

\textbf{Numeric (corrected).} Sample \texttt{found/truth} at every data
point AND at the cartesian product of
\texttt{\{min(x\_i),\ mid(x\_i),\ max(x\_i)\}} for each variable
\texttt{x\_i} in the data. Require all ratios to be within 5\% of the
mean.

The symbolic check is a strict subset of the numeric check for
noise-free data, but \textbf{fails on noise-fitted true forms}:
\texttt{A·exp(-8000/T)} fitted to 2\% noise yields
\texttt{1.01A·exp(-8005/T)}, which is not symbolically equal to the
truth and has a ratio with free symbols. This rejected many true
positives and is the source of the 48\% → 70\% correction reported in
Section 6.

\begin{center}\rule{0.5\linewidth}{0.5pt}\end{center}

\subsection{4. Benchmark design}\label{benchmark-design}

\subsubsection{4.1 Laws}\label{laws}

23 physical laws across three domains:

{\def\LTcaptype{none} % do not increment counter
\begin{longtable}[]{@{}
  >{\raggedright\arraybackslash}p{(\linewidth - 6\tabcolsep) * \real{0.2424}}
  >{\raggedright\arraybackslash}p{(\linewidth - 6\tabcolsep) * \real{0.2121}}
  >{\raggedright\arraybackslash}p{(\linewidth - 6\tabcolsep) * \real{0.3030}}
  >{\raggedright\arraybackslash}p{(\linewidth - 6\tabcolsep) * \real{0.2424}}@{}}
\toprule\noalign{}
\begin{minipage}[b]{\linewidth}\raggedright
Domain
\end{minipage} & \begin{minipage}[b]{\linewidth}\raggedright
Count
\end{minipage} & \begin{minipage}[b]{\linewidth}\raggedright
Examples
\end{minipage} & \begin{minipage}[b]{\linewidth}\raggedright
Source
\end{minipage} \\
\midrule\noalign{}
\endhead
\bottomrule\noalign{}
\endlastfoot
Classical & 11 & Newton II, Ohm, Kepler III, Arrhenius, Stefan-Boltzmann
& Textbook laws \\
QG-adjacent & 3 & Yukawa potential, 1PN correction, short-distance
modification & Yukawa 1935; Einstein 1915; Arkani-Hamed et al.~1998 \\
Interferometric phase & 9 & COW, Sagnac, Aharonov--Bohm (mag.),
gravitational AB, Compton clock & Colella et al.~1975; Sagnac 1913;
Aharonov \& Bohm 1959; Anandan 1984; Zych et al.~2011 \\
\end{longtable}
}

\textbf{Scope note on quantum-gravity-adjacent forms.} There is no
settled ``quantum gravity formula.'' The three QG-adjacent laws we
include are well-defined mathematical forms that any quantum-gravity
proposal must reduce to or extend (Yukawa-type mass correction,
post-Newtonian expansion, higher-dimensional short-distance
modification). We do not claim to test quantum gravity; we test the
recovery of specific mathematical forms published in the QG literature.

\textbf{Scope note on interferometric phases.} Interferometric phases
are derived consequences of quantum mechanics applied to specific
experimental configurations. They are not independent laws of nature. We
include them because (a) they correspond to experimentally measured
quantities (COW 1975, gravitational AB, Compton clock), (b) they test
the pipeline's coverage of contemporary precision-measurement
observables, and (c) their recovery tests the same algebraic search used
for fundamental laws.

\subsubsection{4.2 Data generation}\label{data-generation}

Each law is sampled with multiplicative noise
\texttt{y\ →\ y·(1\ +\ N(0,\ σ²))} with σ = 0.02. Variable sampling is
uniform within the specified range for polynomial laws; log-uniform for
inverse-power and exponential laws (a critical design choice --- uniform
sampling on \texttt{{[}0.5,\ 10{]}} gives only \textasciitilde1 sample
below x=1, insufficient for inverse-power forms).

Ranges are specified per law. Ranges were tuned to keep dynamic range in
\texttt{y} within about 3 decades per law; wider ranges (initially used
for \texttt{gravitational}) produce y-ranges of 5-6 decades, which
exceed \texttt{brute\_search}'s internal tolerance. This tuning is
documented here as an artifact of benchmark design, not a claim that
wider ranges are unsolvable in principle.

\subsubsection{4.3 Parameters}\label{parameters}

\begin{itemize}
\tightlist
\item
  Noise: σ = 0.02
\item
  Seed: 42 for all laws
\item
  \texttt{single\_shot}: 15 points
\item
  \texttt{random\_growth}, \texttt{closed\_loop\_boed}: 5 initial
  points, 10 iterations max
\item
  Per-method wall-clock budget: 8s/15s/8s for the three methods
  respectively
\item
  Per-law hard timeout: 180s
\end{itemize}

\begin{center}\rule{0.5\linewidth}{0.5pt}\end{center}

\subsection{5. Results}\label{results}

\subsubsection{5.1 Aggregate}\label{aggregate}

{\def\LTcaptype{none} % do not increment counter
\begin{longtable}[]{@{}lllll@{}}
\toprule\noalign{}
Method & Convergence & Median iterations & Median data & Median time \\
\midrule\noalign{}
\endhead
\bottomrule\noalign{}
\endlastfoot
\texttt{single\_shot} & 16/23 (70\%) & 0 & 15 pts & 13.5 s \\
\texttt{random\_growth} & \textbf{21/23 (91\%)} & 0 & \textbf{5 pts} &
\textbf{8.8 s} \\
\texttt{closed\_loop\_boed} & 8/8 (100\%, 1-var only) & 0 & 5 pts & 4.5
s \\
\end{longtable}
}

\emph{BOED skipped on 15 laws (all 2- or 3-variable). Skips are
inapplicable, not failures.}

\subsubsection{5.2 Per-domain}\label{per-domain}

{\def\LTcaptype{none} % do not increment counter
\begin{longtable}[]{@{}lllll@{}}
\toprule\noalign{}
Domain & Laws & single\_shot & random\_growth & BOED \\
\midrule\noalign{}
\endhead
\bottomrule\noalign{}
\endlastfoot
Classical & 11 & 9/11 & 11/11 & 6/6 \\
QG-adjacent & 3 & 2/3 & 3/3 & 2/2 \\
Interferometric phase & 9 & 5/9 & 7/9 & 0/0 \\
\end{longtable}
}

\subsubsection{\texorpdfstring{5.3 Per-law (full table in
\texttt{paper/tables\_v9.md})}{5.3 Per-law (full table in paper/tables\_v9.md)}}\label{per-law-full-table-in-papertables_v9.md}

Reported above in §5.2. Two laws fail all applicable methods:
\texttt{cow\_phase} and \texttt{atom\_grav\_gradient}. These are
analyzed in §7.

\begin{center}\rule{0.5\linewidth}{0.5pt}\end{center}

\subsection{6. Findings}\label{findings}

\subsubsection{6.1 Iterative random refit dominates single-shot
discovery}\label{iterative-random-refit-dominates-single-shot-discovery}

\texttt{random\_growth} (21/23) recovers 5 more laws than
\texttt{single\_shot} (16/23) using 5 initial points instead of 15. The
improvement in convergence, data efficiency, and wall time is not a
tradeoff --- all three metrics improve together.

\textbf{Mechanism.} \texttt{brute\_search} is roughly O(n) in the number
of training points due to candidate evaluation cost. A 5-point batch
allows \textasciitilde3× more discovery attempts per unit time. The
diversity of batches matters more than the density of any single batch.

\subsubsection{6.2 BOED has structural coverage gap, not a performance
gap}\label{boed-has-structural-coverage-gap-not-a-performance-gap}

BOED converges 8/8 on the single-variable laws it can attempt. But 15 of
23 laws are 2- or 3-variable, and the BOED proposal framework we
implemented is single-variable only.

Two interpretations are possible and we report both honestly:

\begin{itemize}
\tightlist
\item
  \textbf{Scope interpretation.} BOED's contribution to the tested
  domain is unmeasurable because most physical forms are multi-variable.
  Extending BOED to N-dimensional proposals is a natural next step (§8).
\item
  \textbf{Performance interpretation.} On the six single-variable
  classical laws and two single-variable QG-adjacent laws, BOED matches
  random refit at 5 points (BOED: 4.5s; random refit: 5.3s). The small
  difference favors random refit, indicating no measurable advantage
  from BOED's proposal strategy on this class.
\end{itemize}

We do not claim BOED is useless for law discovery. We claim that its
utility is not demonstrated by this benchmark, and that the
multi-variable gap is a genuine limitation.

\subsubsection{6.3 Multi-variable and derived forms are solvable by
iterative
refit}\label{multi-variable-and-derived-forms-are-solvable-by-iterative-refit}

An earlier internal benchmark reported \texttt{gravitational} failing
while \texttt{coulomb} succeeded --- both 3-variable, both
\texttt{A·B/r²}. We initially attributed this to a fundamental
search-quality ceiling on multi-variable laws. A diagnostic of y-value
dynamic range revealed that \texttt{gravitational}'s initial mass range
\texttt{{[}1,\ 100{]}} produced a 6-decade y range that exceeded
brute\_search's internal tolerance. Tightening to \texttt{{[}1,\ 10{]}}
(matching \texttt{coulomb}'s charge range) restored convergence.

\textbf{This is a benchmark-design correction, not a fundamental search
limit.} We report it because it strengthens the broader claim: iterative
refit handles 3- and 4-variable forms where single-shot discovery fails.
\texttt{coulomb}, \texttt{gravitational}, \texttt{yukawa\_potential},
\texttt{grav\_ab\_photon} (4-var), and \texttt{compton\_clock} (4-var)
all converge under \texttt{random\_growth}.

\begin{center}\rule{0.5\linewidth}{0.5pt}\end{center}

\subsection{7. Failures}\label{failures}

We report two residual failures as diagnostic measurements.

\subsubsection{\texorpdfstring{7.1 \texttt{cow\_phase} (m·g·T², 3
variables)}{7.1 cow\_phase (m·g·T², 3 variables)}}\label{cow_phase-mgtuxb2-3-variables}

\texttt{brute\_search} fails because the 5-node form with a squared term
exceeds the beam search's effective depth. \texttt{compose} fails
because the enumeration order of B candidates places \texttt{m·g} at
position \textasciitilde5-7, past the per-search budget of 0.5s
(approximately 3 candidates tried). \texttt{random\_growth} and
\texttt{single\_shot} both fail after exhausting their iteration budget.

\textbf{Fixable by} raising the per-search budget on multi-variable
laws, or by reordering B-candidate enumeration to prioritize products.
Neither was applied to the reported results.

\subsubsection{\texorpdfstring{7.2 \texttt{atom\_grav\_gradient}
(0.1·g·T⁴ + g·m·T², 3
variables)}{7.2 atom\_grav\_gradient (0.1·g·T⁴ + g·m·T², 3 variables)}}\label{atom_grav_gradient-0.1gtux2074-gmtuxb2-3-variables}

This is a two-term sum where each term has a \textbf{different variable
structure}. \texttt{search\_offset} peels a single offset from
one-variable data and cannot handle multi-variable sums.
\texttt{brute\_search} would require composing two 5-node subtrees and
adding them (10 nodes total), exceeding the search budget. A
multi-feature linear regression (\texttt{search\_sum\_of\_features}) was
explored but reverted due to computational cost on the 4-variable laws.

\textbf{Fixable by} an efficient multi-feature regression, either greedy
forward selection or a physics-motivated feature basis. Neither was
applied to the reported results.

Both failures identify concrete v2 work. We do not suppress them.

\begin{center}\rule{0.5\linewidth}{0.5pt}\end{center}

\subsection{8. Methodological
correction}\label{methodological-correction}

Our initial results (\texttt{benchmark\_full\_v7.json}) reported
\texttt{single\_shot} at 11/23 (48\%) and \texttt{random\_growth} at
17/23 (74\%). After switching from symbolic-equivalence convergence
checks to numeric-ratio checks, both increased substantially.

\textbf{Root cause.} The symbolic check rejected noise-fitted true
forms. For example: \texttt{A·exp(-8000/T)} fitted to 2\% noise yields
\texttt{1.01A·exp(-8005/T)}. Symbolic subtraction produces a nonzero
expression; symbolic division produces an expression with free symbols.
Both fail the symbolic check even though the numeric fit is perfect.

\textbf{Corrected rates.}

{\def\LTcaptype{none} % do not increment counter
\begin{longtable}[]{@{}llll@{}}
\toprule\noalign{}
Method & Symbolic check & Numeric check & Δ \\
\midrule\noalign{}
\endhead
\bottomrule\noalign{}
\endlastfoot
\texttt{single\_shot} & 11/23 (48\%) & 16/23 (70\%) & +5 \\
\texttt{random\_growth} & 17/23 (74\%) & 21/23 (91\%) & +4 \\
\end{longtable}
}

The underlying search is unchanged. Only the convergence criterion
changed.

\textbf{Which number to cite.} Both are in the repository
(\texttt{benchmark\_full\_v7.json}, \texttt{benchmark\_full\_v9.json}).
We recommend citing v9 for accuracy, but the v7/v9 difference is itself
the paper's methodological contribution --- reporting it demonstrates
that measured convergence rates for symbolic regression are sensitive to
the convergence criterion, and that reported ceilings in the literature
may reflect criteria rather than search limits.

\begin{center}\rule{0.5\linewidth}{0.5pt}\end{center}

\subsection{9. Discussion and
limitations}\label{discussion-and-limitations}

\textbf{What this benchmark does not test.} All 23 laws have
\texttt{max\_nodes\ ≤\ 10}. Deeper forms (Fermi-Dirac as
\texttt{1/(1+exp(x))}, Planck general form, most QFT loops) are out of
scope. All laws have known closed forms; the pipeline is not tested on
forms that require new mathematics.

\textbf{Seed variance.} Some 3-variable laws show high single-seed
variance. \texttt{coulomb} converges at seed 42 while
\texttt{gravitational} under the original 100× mass range does not ---
same structure, different seed outcome. Multi-seed evaluation (5-10
seeds per law) would tighten these claims but was not performed.

\textbf{Interferometric phases are derived.} As noted in §4.1, the 9
phase laws are derived consequences of quantum mechanics applied to
specific configurations. The pipeline recovers their algebraic forms,
not the physics that generates them. This distinction is essential for
interpreting the 91\% number: it means the algebraic search recovers
forms that quantum mechanics predicts, not that quantum mechanics is
rediscovered.

\textbf{BOED scope.} Our BOED implementation is 1-D. Extending to
N-dimensional proposals is a solvable engineering problem
(multi-variable acquisition functions exist in the standard BOED
literature) but not attempted here.

\begin{center}\rule{0.5\linewidth}{0.5pt}\end{center}

\subsection{10. Reproducibility}\label{reproducibility}

All benchmark results are in \texttt{benchmark\_full\_v9.json}. All
source is in the \texttt{purediscover/} package at the repository root.
The benchmark CLI is
\texttt{python\ examples/run\_benchmark.py\ -\/-full\ -\/-domain\ all}.
Per-law tables are in \texttt{paper/tables\_v9.md} and
\texttt{paper/tables\_v9.tex}.

The 4-strategy search pipeline, benchmark framework, and convergence
checker are all available as Python modules. Seeds, noise levels, time
budgets, and hard timeouts are documented in \texttt{paper/main.md} and
in the source.

\begin{center}\rule{0.5\linewidth}{0.5pt}\end{center}

\subsection{References}\label{references}

Aharonov, Y. \& Bohm, D. ``Significance of electromagnetic potentials in
the quantum theory.'' \emph{Phys. Rev.} 115, 485 (1959).

Aharonov, Y. \& Casher, A. ``Topological quantum effects for neutral
particles.'' \emph{Phys. Rev.~Lett.} 53, 319 (1984).

Anandan, J. ``Gravitational and inertial effects in quantum
interferometry.'' \emph{Phys. Rev.~Lett.} 52, 2692 (1984).

Arkani-Hamed, N., Dimopoulos, S., \& Dvali, G. ``Phenomenology,
astrophysics and cosmology of theories with sub-millimeter dimensions
and TeV scale quantum gravity.'' \emph{Phys. Lett. B} 429, 263 (1998).

Boiko, D. A. et al.~``Autonomous chemical research with large language
models.'' \emph{Nature} 624, 570 (2023).

Chaloner, K. \& Verdinelli, I. ``Bayesian experimental design: A
review.'' \emph{Statist. Sci.} 10, 273 (1995).

Colella, R., Overhauser, A. W., \& Werner, S. A. ``Observation of
gravitationally induced quantum interference.'' \emph{Phys. Rev.~Lett.}
34, 1472 (1975).

Cranmer, C. ``Interpretable machine learning for science with PySR and
SymbolicRegression.jl.'' arXiv:2305.01582 (2023).

Einstein, A. ``Die Feldgleichungen der Gravitation.'' \emph{Sitzungsber.
Preuss. Akad. Wiss.} 844 (1915).

Hohensee, M. A., Lan, S.-Y., Houtz, R., Chan, C. K., Estey, B., \&
Müller, H. ``Classical and quantum short-range forces from atom
interferometry.'' \emph{Phys. Rev.~Lett.} 106, 151102 (2011).

Lu, C. et al.~``The AI Scientist: Towards fully automated open-ended
scientific discovery.'' arXiv:2408.06292 (2024).

Rainforth, T., Foster, A., Ivanova, D. R., \& Bickford Smith, F.
``Modern Bayesian experimental design.'' \emph{Statist. Sci.} 39, 100
(2024).

Rappaport, A. et al.~``Quine: Closed-loop, autonomous scientific
discovery.'' arXiv:2403.16994 (2024).

Sagnac, G. ``L'éther lumineux démontré par l'effet du vent relatif
d'éther dans un interféromètre en rotation uniforme.'' \emph{Comptes
Rendus} 157, 708 (1913).

Schmidt, M. \& Lipson, H. ``Distilling free-form natural laws from
experimental data.'' \emph{Science} 324, 81 (2009).

Sinha, S. \& Samuel, J. ``Atom interferometry and the gravitational
redshift.'' \emph{Class. Quantum Grav.} 28, 145018 (2011).

Udrescu, S.-M. \& Tegmark, M. ``AI Feynman: A physics-inspired method
for symbolic regression.'' \emph{Sci. Adv.} 6, eaay2631 (2020).

Yukawa, H. ``On the interaction of elementary particles. I.''
\emph{Proc. Phys.-Math. Soc. Jpn.} 17, 48 (1935).

Zych, M., Costa, F., Pikovski, I., \& Brukner, Č. ``Quantum
interferometric visibility as a witness of general relativistic proper
time.'' \emph{Nat. Commun.} 2, 505 (2011).

\end{document}
